\chapter{Adrenocorticotropic Hormone (ACTH)}

\section*{What Is This Test?}
Adrenocorticotropic hormone (ACTH, also called corticotropin) is a 39-amino-acid polypeptide hormone produced by corticotrope cells of the anterior pituitary gland. ACTH is synthesized from a large precursor molecule, pro-opiomelanocortin (POMC), which is cleaved by prohormone convertase 1 (PC1) to produce ACTH and beta-lipotropin. ACTH stimulates the adrenal cortex (primarily zona fasciculata and zona reticularis) to produce cortisol and adrenal androgens (DHEA, androstenedione). ACTH secretion is regulated by corticotropin-releasing hormone (CRH) from the hypothalamus and arginine vasopressin (AVP), and is under negative feedback control by cortisol. ACTH exhibits a circadian rhythm similar to cortisol: peak levels at 6--8 AM and nadir at 11 PM--1 AM. The plasma half-life of ACTH is only 7--12 minutes, reflecting rapid clearance. The ACTH test is primarily used in conjunction with cortisol measurement to determine the etiology of adrenal disorders: elevated ACTH with low cortisol indicates primary adrenal insufficiency \cite{bornstein2016diagnosis}; elevated ACTH with elevated cortisol indicates ACTH-dependent Cushing syndrome \cite{lacroix2015cushing}; low ACTH with elevated cortisol indicates ACTH-independent Cushing syndrome \cite{lacroix2015cushing}; low ACTH with low cortisol indicates secondary adrenal insufficiency \cite{bancos2015diagnosis}. ACTH is measured in plasma (not serum) because it is susceptible to proteolytic degradation and adsorption to glass tubes.

\section*{Alternative Names}
Corticotropin; Adrenocorticotropin; Plasma ACTH; Serum ACTH; Cosyntropin (synthetic ACTH 1-24, trade name Cortrosyn); ACTH 1-39

\section*{Principle}
ACTH is measured using two-site sandwich immunometric (chemiluminescent) assays. Two monoclonal antibodies directed against different epitopes on the ACTH molecule (typically the N-terminal and C-terminal regions) are used: a capture antibody on a solid phase and a detection antibody labeled with a chemiluminescent tag. The assay specifically detects intact ACTH (1--39). Common platforms: Roche Cobas ECLIA, Siemens IMMULITE/Atellica, and Diasorin Liaison. The functional sensitivity is typically <5 pg/mL for high-sensitivity assays. The major pre-analytical challenge for ACTH measurement is its instability: ACTH is rapidly degraded by plasma proteases at room temperature (half-life in vitro is 1--2 hours at 25 degrees C). For this reason, blood must be collected into a pre-chilled EDTA tube (purple top), placed immediately on ice, centrifuged at 4 degrees C within 1 hour, and the plasma should be aliquoted, frozen at -20 degrees C or colder, and transported on dry ice. Aprotinin (a protease inhibitor) may be added to the collection tube to prevent degradation. EDTA plasma is required (not serum or heparin) because ACTH adsorbs to glass, so plastic tubes only. LC-MS/MS is not used for ACTH measurement because of the antibody requirement for specificity and the low circulating concentrations.

\section*{History}
ACTH was first isolated from pituitary extracts in the 1930s. The first ACTH bioassay (adrenal ascorbic acid depletion in hypophysectomized rats) was developed by Sayers and colleagues in 1948. The radioimmunoassay for ACTH was developed by Yalow and Berson in 1968 (the same team that first developed the RIA for insulin). The POMC precursor was discovered in 1979 by Nakanishi and colleagues, and Eipper and Mains, revealing that ACTH is one of several peptides produced from the POMC prohormone (also including beta-endorphin, MSH, and CLIP). The development of the CRH stimulation test for Cushing syndrome by Chrousos, Nieman, and colleagues in the 1980s, and the inferior petrosal sinus sampling technique by Oldfield and Doppman at the NIH, revolutionized the differential diagnosis of ACTH-dependent Cushing syndrome \cite{oldfield1991petrosal}. The Endocrine Society guidelines for Cushing syndrome (2008, updated 2015) and adrenal insufficiency (2016) codified the use of ACTH in diagnostic algorithms \cite{nieman2015treatment, bornstein2016diagnosis}. Current challenges include assay harmonization, since ACTH immunoassays vary across platforms (inter-assay variation 15--25\%) \cite{findling2017diagnosis}.

\section*{Why This Test Is Ordered}
\begin{itemize}
  \item Classifying adrenal insufficiency as primary (elevated ACTH, low cortisol) vs. secondary (low or inappropriately normal ACTH, low cortisol) vs. tertiary (hypothalamic, low ACTH, low cortisol)
  \item Classifying Cushing syndrome as ACTH-dependent (Cushing disease or ectopic ACTH, elevated or inappropriately normal ACTH) vs. ACTH-independent (adrenal adenoma or carcinoma, suppressed ACTH <5--10 pg/mL)
  \item Differentiating Cushing disease (pituitary, moderate ACTH elevation) from ectopic ACTH syndrome (markedly elevated ACTH, often >200 pg/mL, with severe hypokalemia)
  \item Investigating the source of ACTH in ACTH-dependent Cushing syndrome via inferior petrosal sinus sampling with ACTH measurement
  \item Monitoring Nelson syndrome (post-bilateral adrenalectomy ACTH-secreting pituitary adenoma growth) --- rising ACTH >500 pg/mL with skin hyperpigmentation
  \item Evaluation of congenital adrenal hyperplasia (elevated ACTH due to impaired cortisol synthesis) --- monitoring adequacy of glucocorticoid therapy
  \item Diagnosis of ectopic ACTH-producing tumors (small cell lung cancer, bronchial carcinoid, thymic carcinoid, medullary thyroid carcinoma, pancreatic NETs)
  \item Investigating hyperpigmentation --- elevated ACTH and MSH (from POMC) in primary adrenal insufficiency causes characteristic hyperpigmentation of sun-exposed areas, palmar creases, buccal mucosa, and scars
\end{itemize}

\section*{Primary vs Secondary Test}
This is a secondary classification test. ACTH is always interpreted in conjunction with a simultaneously drawn cortisol level. It is never used as a screening test for adrenal disorders. ACTH measurement is indicated when cortisol results are abnormal (low or elevated), to determine whether the problem is adrenal (ACTH-independent) or pituitary/hypothalamic/ectopic (ACTH-dependent). The Endocrine Society guidelines specify ACTH measurement as the first step after confirming abnormal cortisol (elevated or insufficient).

\section*{How This Test Is Performed}
A healthcare professional draws a blood sample from a vein into a pre-chilled EDTA (purple-top) plastic tube. The tube MUST be placed on ice immediately after collection. The specimen is centrifuged at 4 degrees C within 1 hour of collection; plasma is separated, transferred to a plastic transport tube, and frozen at -20 degrees C or colder. The test must be drawn simultaneously with serum cortisol (gold-top tube) at 8--9 AM to capture the peak ACTH-cortisol relationship. For inferior petrosal sinus sampling (IPSS), catheters are placed in both inferior petrosal sinuses via femoral vein access by interventional radiology. ACTH is measured simultaneously from both petrosal sinuses and a peripheral vein at baseline and 3, 5, 10, and 15 minutes after CRH administration (1 mcg/kg IV). A central-to-peripheral ACTH gradient >2:1 basal and >3:1 post-CRH confirms pituitary Cushing disease. ACTH results must be correlated with the simultaneously drawn cortisol level. Results from improperly handled specimens (warm, delayed) are invalid due to ACTH degradation.

\section*{What Preparation Is Needed}
The patient should fast for 8--10 hours before testing and be in a relaxed, non-stressed state (stress and pain elevate ACTH). Morning samples (8--9 AM) are drawn for basal ACTH and cortisol because this represents the circadian peak and is when reference ranges apply. Afternoon or midnight levels are used in specialized testing (e.g., midnight ACTH in CRH testing, or IPSS). For IPSS, the patient must be fasting overnight and should be sedated to minimize stress-induced ACTH release. High-dose biotin (>5 mg/day) must be discontinued 48 hours before testing. Drugs affecting ACTH: glucocorticoids suppress ACTH (patients on glucocorticoid therapy will have suppressed ACTH); CRH and vasopressin stimulate ACTH; amphetamines and metoclopramide stimulate ACTH; opioid analgesics and GABA agonists may suppress ACTH. The test must NOT be drawn after glucocorticoid administration unless specifically indicated (e.g., cortisol after dexamethasone suppression).

\section*{Contraindications and Precautions}
Pre-analytical factors are critical: incorrect specimen handling is the most common cause of falsely low ACTH results. Blood drawn into room-temperature tubes, delayed centrifugation (>1 hour), or failure to freeze immediately results in ACTH proteolytic degradation, producing falsely low values that can mislead diagnosis. Hemolysis can falsely elevate ACTH in some assays because of interference with chemiluminescent detection. Icterus and lipemia also interfere. Heterophile antibodies and HAMA can cause false elevations. Patients on glucocorticoid therapy will have suppressed ACTH (expected) --- this is not a false result but reflects HPA axis suppression. ACTH values in the "normal" range must be interpreted in context of cortisol: an ACTH of 15 pg/mL with a cortisol of 3 mcg/dL is inappropriately low (suggesting secondary adrenal insufficiency); an ACTH of 15 pg/mL with a cortisol of 35 mcg/dL is inappropriately normal in the context of hypercortisolism (suggesting Cushing disease). Circadian factors: ACTH peak is 6--8 AM; nadir is 11 PM--1 AM. Loss of diurnal variation (elevated midnight ACTH) is an early feature of Cushing syndrome. ACTH is pulsatile (15--18 pulses per 24 hours), so a single random value may not be representative. The IPSS procedure has risks including: groin hematoma, deep vein thrombosis, pulmonary embolism, and rare brainstem stroke/death (0.2--1\%), and should only be performed at experienced centers.

\section*{Risks}
From venous blood draw: slight pain, bruising, or hematoma at the puncture site. From IPSS: groin hematoma (common), deep venous thrombosis (rare), contrast reaction (rare), transient cranial nerve palsies, and very rarely brainstem infarction or intracranial hemorrhage (reported at 0.2--1\%). IPSS should only be performed at experienced centers with proper monitoring.

\section*{What Happens After the Test}
Results are available within 1--3 days (ACTH is often a send-out test for pre-analytical processing reasons; some reference laboratories perform the test). Interpretation is always in the context of the simultaneously drawn cortisol: (1) Low cortisol + elevated ACTH: primary adrenal insufficiency --- additional testing includes anti-21-hydroxylase antibodies, adrenal CT, and possibly cosyntropin stimulation. (2) Low cortisol + inappropriately normal or low ACTH: secondary/tertiary adrenal insufficiency --- evaluate with pituitary MRI, assess for prior glucocorticoid use, and consider insulin tolerance test if cosyntropin test is equivocal. (3) Elevated cortisol + elevated/inappropriately normal ACTH: ACTH-dependent Cushing syndrome --- requires further testing: high-dose dexamethasone suppression test, CRH stimulation test, pituitary MRI, and IPSS if source is unclear. ACTH >200 pg/mL with hypokalemia strongly suggests ectopic ACTH. (4) Elevated cortisol + suppressed ACTH (<5--10 pg/mL): ACTH-independent Cushing syndrome --- adrenal CT or MRI to identify adenoma or carcinoma.

\section*{Understanding Results}

\subsection*{Below Normal}
\begin{itemize}
  \item \textbf{Suppressed ACTH (<5--10 pg/mL) with elevated cortisol:} ACTH-independent Cushing syndrome. Most commonly an adrenal adenoma (>2 cm, lipid-rich on CT) or adrenal carcinoma (>6 cm, heterogeneous, may secrete androgens causing virilization). ACTH is suppressed by negative feedback from autonomous cortisol production. Less commonly: primary pigmented nodular adrenocortical disease (PPNAD, part of Carney complex, ages 10--30), McCune-Albright syndrome, or bilateral macronodular adrenal hyperplasia (BMAH, older adults).
  \item \textbf{Suppressed ACTH with normal or mildly decreased cortisol:} Exogenous glucocorticoid use --- the most common cause of low ACTH in clinical practice. Any patient taking prednisone-equivalent >5 mg/day for >3 weeks may have HPA axis suppression.
  \item \textbf{Low or inappropriately normal ACTH with low cortisol:} Secondary adrenal insufficiency. The ACTH is inappropriately low relative to the low cortisol. Causes: pituitary adenomas (non-functioning macroadenoma compressing corticotropes), pituitary surgery or radiation, Sheehan syndrome (postpartum pituitary necrosis), traumatic brain injury, lymphocytic hypophysitis, immune checkpoint inhibitor-related hypophysitis (particularly with ipilimumab and combination therapy), isolated ACTH deficiency (rare congenital or acquired).
  \item \textbf{Low ACTH after successful pituitary adenoma resection for Cushing disease:} Post-operative ACTH should be undetectable or very low; cortisol will also be low, confirming surgical remission. Patients require glucocorticoid replacement until HPA axis recovers (typically 6--18 months).
\end{itemize}

\subsection*{Above Normal}
\begin{itemize}
  \item \textbf{Elevated ACTH (>100--300 pg/mL) with elevated cortisol:} ACTH-dependent Cushing syndrome. If the ACTH is moderately elevated (50--150 pg/mL), Cushing disease (pituitary microadenoma) is likely. If ACTH is markedly elevated (>200 pg/mL often >500 pg/mL), ectopic ACTH syndrome is more likely, especially with severe hypokalemia (<3.0 mmol/L). The high-dose dexamethasone suppression test and CRH stimulation test help differentiate (pituitary tumors suppress and respond; ectopic sources typically do neither).
  \item \textbf{Cushing disease (pituitary ACTH adenoma):} ACTH typically 50--150 pg/mL. Cortisol suppresses >50\% from baseline on high-dose dexamethasone. ACTH and cortisol rise >35--50\% and >20\% respectively after CRH. Pituitary adenoma visible on MRI in 50--60\%. IPSS confirms pituitary source.
  \item \textbf{Ectopic ACTH syndrome:} ACTH often >200--500 pg/mL, may exceed 1000 pg/mL. Associated with severe hypercortisolism, resistant hypokalemia, and metabolic alkalosis. No suppression on high-dose dexamethasone and no response to CRH. Most common tumors: small cell lung carcinoma, bronchial carcinoid (often occult, small, slow-growing), thymic carcinoid, pancreatic neuroendocrine tumors, medullary thyroid carcinoma, pheochromocytoma. These tumors process POMC abnormally, producing high molecular weight ACTH precursors.
  \item \textbf{Elevated ACTH (>100 pg/mL, often >500--1000 pg/mL) with low cortisol:} Primary adrenal insufficiency (Addison disease). The absence of cortisol negative feedback drives unopposed CRH and ACTH secretion. Elevated ACTH and POMC-derived peptides (including MSH) cause hyperpigmentation, a pathognomonic sign. Causes: autoimmune adrenalitis (80\% in developed countries, anti-21-hydroxylase antibodies positive), infectious (tuberculosis, HIV, fungal), adrenal hemorrhage (Waterhouse-Friderichsen syndrome), congenital adrenal hyperplasia, adrenoleukodystrophy, medications (ketoconazole, etomidate, mitotane, checkpoint inhibitors).
  \item \textbf{Elevated ACTH in congenital adrenal hyperplasia:} 21-hydroxylase deficiency (90--95\% of CAH) causes impaired cortisol synthesis, loss of negative feedback, and elevated ACTH driving adrenal hyperplasia and excess androgen production. 17-hydroxyprogesterone is elevated. Other CAH forms (11-beta-hydroxylase, 3-beta-hydroxysteroid dehydrogenase deficiencies) also elevate ACTH.
  \item \textbf{Nelson syndrome:} Rising ACTH (often >500--1000 pg/mL) with cutaneous hyperpigmentation in patients who underwent bilateral adrenalectomy for Cushing disease. Occurs in 8--38\% of bilaterally adrenalectomized patients. Caused by growth of pre-existing pituitary microadenoma freed from cortisol negative feedback. Requires pituitary MRI and may need surgery or radiation.
\end{itemize}

\section*{Normal Results}
\begin{itemize}
  \item \textbf{Plasma ACTH (8 AM, supine):} 10--60 pg/mL (2--13 pmol/L) --- varies by assay; typical range 7--63 pg/mL (Roche, Siemens)
  \item \textbf{Plasma ACTH (4 PM):} 5--30 pg/mL
  \item \textbf{Plasma ACTH (midnight, sleeping):} <5--10 pg/mL
  \item \textbf{ACTH in Cushing disease:} 50--150 pg/mL (moderately elevated or "inappropriately normal" in context of hypercortisolism)
  \item \textbf{ACTH in ectopic ACTH syndrome:} >200 pg/mL, often >500 pg/mL (markedly elevated)
  \item \textbf{ACTH in primary adrenal insufficiency:} >100 pg/mL, often >500 pg/mL (markedly elevated with absent cortisol feedback)
  \item \textbf{ACTH in ACTH-independent Cushing (adrenal):} <5--10 pg/mL (suppressed)
  \item \textbf{ACTH in secondary adrenal insufficiency:} <20 pg/mL or inappropriately "normal" with low cortisol
  \item \textbf{Note:} ACTH assays vary significantly between manufacturers. Always use the same assay for serial monitoring. Always interpret ACTH simultaneously with cortisol. The specimen MUST be collected and handled properly (pre-chilled EDTA tube, on ice, centrifuged at 4 degrees C, frozen within 1 hour).
\end{itemize}

\section*{Follow-up Tests}
\begin{itemize}
  \item \textbf{Simultaneous serum cortisol:} ACTH is NEVER interpreted in isolation. The ACTH-cortisol pair is fundamental to classifying adrenal disorders. The 8 AM simultaneous draw is standard.
  \item \textbf{High-dose dexamethasone suppression test (8 mg overnight):} Dexamethasone 8 mg at 11 PM, measure cortisol 8 AM next day. >50\% suppression of cortisol suggests pituitary Cushing disease. Lack of suppression suggests ectopic ACTH or adrenal source.
  \item \textbf{CRH stimulation test:} Ovine CRH 1 mcg/kg IV. ACTH increase >35--50\% at 15--30 minutes suggests Cushing disease. Ectopic ACTH-producing tumors usually do not respond to CRH.
  \item \textbf{Inferior petrosal sinus sampling (IPSS) with CRH:} Gold standard for confirming pituitary Cushing disease. Central-to-peripheral ACTH gradient >2:1 basal and >3:1 post-CRH. Also lateralizes the adenoma within the pituitary (inter-sinus gradient >1.4:1).
  \item \textbf{Pituitary MRI with gadolinium contrast:} Thin-section (1--3 mm) coronal and sagittal T1-weighted images before and after contrast. Detects pituitary adenoma in 50--60\% of Cushing disease. Tumor <6 mm may not be visible.
  \item \textbf{Adrenal CT (non-contrast, with washout):} For ACTH-independent Cushing syndrome (suppressed ACTH). Assesses adrenal adenoma characteristics (size, lipid content, washout).
  \item \textbf{High-resolution CT chest/abdomen/pelvis:} Search for ectopic ACTH source (occult bronchial carcinoid most common). Small tumors (<1--2 cm) may be missed on conventional CT. Octreotide scan or Ga-68 DOTATATE PET/CT helps detect somatostatin receptor-positive neuroendocrine tumors.
  \item \textbf{Anti-21-hydroxylase antibodies:} Confirm autoimmune Addison disease in primary adrenal insufficiency. Positive in >80\%.
  \item \textbf{Very long-chain fatty acids (VLCFA):} Screen for X-linked adrenoleukodystrophy in males with primary adrenal insufficiency (elevated VLCFA, especially C26:0).
  \item \textbf{17-hydroxyprogesterone:} Screen for 21-hydroxylase deficiency (congenital adrenal hyperplasia) in patients with elevated ACTH and low cortisol.
\end{itemize}

\section*{References}
\bibliographystyle{unsrtnat}
\bibliography{references}

\clearpage
\section*{Regulatory Status and Authorization Date}
Regulatory status applies to the specific assay, instrument, manufacturer, and intended use rather than to the test name alone. No single approval date can be assigned to this test category. For a commercial assay, verify the current product in the relevant regulator's database: the U.S. Food and Drug Administration (FDA) clearance, approval, or authorization; India's Central Drugs Standard Control Organisation (CDSCO) licence or permission; the European Union In Vitro Diagnostic Regulation (EU IVDR) CE certificate issued through the applicable conformity-assessment route; or the United Kingdom Medicines and Healthcare products Regulatory Agency (MHRA) registration and UKCA/CE status. Laboratory-developed tests may not have an individual FDA, CDSCO, EU IVDR, or MHRA approval date.
\clearpage
